\section{Sujet n\textsuperscript{o}\,28}
\noindent\begin{minipage}{0.5\linewidth}\footnotesize Office du Baccalauréat du Cameroun\end{minipage}%
\hfill\begin{minipage}{0.45\linewidth}\footnotesize\raggedleft Examen : \textbf{Baccalauréat} · Série : \textbf{C/E}\\
Épreuve : \textbf{Mathématiques} · Durée : \textbf{4\,h} · Coef. : \textbf{7}\end{minipage}
\par\smallskip{\color{onbo}\rule{\linewidth}{1pt}}

\subsection*{Partie A — Évaluation des ressources \hfill (15 points)}

\noindent\textbf{Exercice 1.} \hfill\textit{(3 points)}\\
Au marché à bétail de Maroua, les bovins proviennent de trois éleveurs : $50\,\%$ de
l'éleveur~$M$, $30\,\%$ de l'éleveur~$N$ et $20\,\%$ de l'éleveur~$P$. Un vétérinaire
estime que la proportion de bêtes atteintes d'une parasitose est de $4\,\%$ chez~$M$,
$10\,\%$ chez~$N$ et $15\,\%$ chez~$P$. On choisit au hasard un bovin du marché. On note
$M$, $N$, $P$ les événements « le bovin provient de l'éleveur correspondant » et $S$
l'événement « le bovin est atteint de la parasitose ».
\begin{enumerate}[label=\arabic*)]
\item Construire un arbre pondéré décrivant cette situation. \hfill\textit{(0,75 pt)}
\item Calculer la probabilité que le bovin choisi provienne de l'éleveur~$N$ et soit
atteint de la parasitose. \hfill\textit{(0,5 pt)}
\item Montrer que la probabilité $P(S)$ qu'un bovin du marché soit atteint vaut $0{,}08$. \hfill\textit{(0,75 pt)}
\item Le bovin choisi est atteint de la parasitose. Calculer la probabilité qu'il
provienne de l'éleveur~$P$ (arrondir à $10^{-3}$). \hfill\textit{(1 pt)}
\end{enumerate}

\smallskip
\noindent\textbf{Exercice 2.} \hfill\textit{(3 points)}\\
Un éleveur démarre son cheptel avec $u_0=200$ bovins. Chaque année, le troupeau perd
$10\,\%$ de son effectif (ventes et mortalité) puis l'éleveur achète $30$ nouvelles
bêtes. On note $u_n$ l'effectif du cheptel après $n$ années, de sorte que
$u_{n+1}=0{,}9\,u_n+30$ pour tout $n\in\N$.
\begin{enumerate}[label=\arabic*)]
\item Calculer $u_1$ et $u_2$. \hfill\textit{(0,5 pt)}
\item On pose $v_n=u_n-300$. Montrer que $(v_n)$ est une suite géométrique dont on
précisera la raison et le premier terme. \hfill\textit{(1 pt)}
\item Exprimer $v_n$ puis $u_n$ en fonction de $n$. \hfill\textit{(0,75 pt)}
\item Étudier la convergence de $(u_n)$ et interpréter la limite obtenue. \hfill\textit{(0,75 pt)}
\end{enumerate}

\smallskip
\noindent\textbf{Exercice 3.} \hfill\textit{(4 points)}\\
La taille (en milliers de têtes) d'un cheptel régional est modélisée par une fonction~$y$
du temps~$t$ (en années, $t\ge 0$), solution de l'équation différentielle
$(E):\ y'=0{,}2\,y$ avec la condition initiale $y(0)=5$.
\begin{enumerate}[label=\arabic*)]
\item Résoudre l'équation différentielle $(E)$ et déterminer la solution~$y$ vérifiant
$y(0)=5$. \hfill\textit{(1 pt)}
\item Soit $g$ la fonction définie sur $[0,+\infty[$ par $g(t)=5\,\mathrm{e}^{0{,}2t}$.
Étudier le sens de variation de~$g$ et déterminer sa limite en $+\infty$. \hfill\textit{(1 pt)}
\item Déterminer, à l'aide du modèle, l'année à partir de laquelle le cheptel dépassera
$15\,000$ têtes (arrondir à l'entier). \hfill\textit{(1 pt)}
\item Calculer la valeur moyenne $\dfrac{1}{5}\displaystyle\int_0^{5} g(t)\,\dd t$ du
cheptel sur les cinq premières années (arrondir à $10^{-2}$). \hfill\textit{(1 pt)}
\end{enumerate}

\smallskip
\noindent\textbf{Exercice 4.} \hfill\textit{(5 points)}\\
Le plan complexe est muni d'un repère orthonormé direct $(O\,;\,\vec u,\vec v)$. Sur un
plan cadastral du marché, on repère trois bornes $A$, $B$ et $C$ d'affixes respectives
$z_A=2+2i$, $z_B=6+2i$ et $z_C=4+6i$.
\begin{enumerate}[label=\arabic*)]
\item Calculer $\dfrac{z_C-z_A}{z_B-z_A}$ ; en déduire la nature du triangle $ABC$. \hfill\textit{(1,25 pt)}
\item On considère la similitude directe $S$ qui transforme $A$ en $B$ et $B$ en $C$.
Déterminer l'écriture complexe de $S$ sous la forme $z'=az+b$. \hfill\textit{(1,5 pt)}
\item Déterminer le rapport et l'angle de $S$. \hfill\textit{(0,75 pt)}
\item Déterminer l'affixe~$\omega$ du centre~$\Omega$ de la similitude~$S$. \hfill\textit{(1,5 pt)}
\end{enumerate}

\subsection*{Partie B — Évaluation des compétences \hfill (5 points)}

\noindent\textit{Monsieur Bakari, éleveur au marché à bétail de Maroua, te sollicite, en
tant qu'élève de Terminale~C, pour trois décisions concernant son exploitation. Il dispose
de \SI{120}{\metre} de barrière métallique et veut clôturer un enclos rectangulaire
\emph{adossé à un long hangar} déjà construit (le côté longeant le hangar n'est pas
clôturé). Par ailleurs, son cheptel actuel de $400$ bovins croît, selon un modèle continu,
au taux annuel de \SI{12}{\percent} : il est modélisé par $N(t)=400\,\mathrm{e}^{0{,}12t}$,
où $t$ désigne le nombre d'années. Enfin, à la vaccination, \SI{6}{\percent} des bêtes
réagissent mal au vaccin ; un contrôleur prélève $5$ bêtes au hasard dans un grand
troupeau.}

\begin{enumerate}[label=\textbf{Tâche \arabic*.},leftmargin=2.4em]
\item Déterminer les dimensions de l'enclos rectangulaire d'\emph{aire maximale} qu'il
peut clôturer, ainsi que cette aire. \hfill\textit{(1,5 pt)}
\item Déterminer le nombre d'années (entier) au bout duquel le cheptel aura doublé. \hfill\textit{(1,5 pt)}
\item Déterminer la probabilité qu'\emph{au moins une} des $5$ bêtes prélevées réagisse
mal au vaccin (arrondir à $10^{-3}$). \hfill\textit{(1,5 pt)}
\end{enumerate}
\par\smallskip{\footnotesize\itshape Présentation générale : 0,5 point.}

% ════════════════════════════════════════════════════
\corrige{du sujet 28}

\noindent\textbf{Partie A — Exercice 1.}
1) Arbre : premier niveau $P(M)=0{,}5$, $P(N)=0{,}3$, $P(P)=0{,}2$ ; second niveau
$P_M(S)=0{,}04$, $P_N(S)=0{,}10$, $P_P(S)=0{,}15$ (et les complémentaires).
2) $P(N\cap S)=P(N)\,P_N(S)=0{,}3\times0{,}10=0{,}03$.
3) Formule des probabilités totales :
$P(S)=0{,}5\times0{,}04+0{,}3\times0{,}10+0{,}2\times0{,}15=0{,}02+0{,}03+0{,}03=0{,}08$.
4) $P_S(P)=\dfrac{P(P\cap S)}{P(S)}=\dfrac{0{,}2\times0{,}15}{0{,}08}=\dfrac{0{,}03}{0{,}08}=0{,}375$.

\noindent\textbf{Exercice 2.}
1) $u_1=0{,}9\times200+30=210$ ; $u_2=0{,}9\times210+30=219$.
2) $v_{n+1}=u_{n+1}-300=0{,}9u_n+30-300=0{,}9u_n-270=0{,}9(u_n-300)=0{,}9\,v_n$ :
$(v_n)$ est géométrique de raison $0{,}9$ et $v_0=u_0-300=-100$.
3) $v_n=-100\times0{,}9^{\,n}$, donc $u_n=300-100\times0{,}9^{\,n}$.
4) Comme $0<0{,}9<1$, $0{,}9^{\,n}\to0$, donc $u_n\to300$ : à long terme le cheptel se
stabilise autour de $300$ bovins.

\noindent\textbf{Exercice 3.}
1) Les solutions de $y'=0{,}2y$ sont $y(t)=K\mathrm{e}^{0{,}2t}$, $K\in\R$. La condition
$y(0)=5$ donne $K=5$ : $y(t)=5\mathrm{e}^{0{,}2t}$.
2) $g'(t)=5\times0{,}2\,\mathrm{e}^{0{,}2t}=\mathrm{e}^{0{,}2t}>0$ : $g$ est strictement
croissante sur $[0,+\infty[$, et $\lim_{t\to+\infty}g(t)=+\infty$.
3) $g(t)>15 \iff \mathrm{e}^{0{,}2t}>3 \iff t>\dfrac{\ln 3}{0{,}2}\approx5{,}49$ : c'est à
partir de la \textbf{6\textsuperscript{e} année} ($g(6)=5\mathrm{e}^{1{,}2}\approx16{,}6$ milliers).
4) $\int_0^5 5\mathrm{e}^{0{,}2t}\dd t=5\big[\tfrac{1}{0{,}2}\mathrm{e}^{0{,}2t}\big]_0^5
=25(\mathrm{e}^{1}-1)=25(2{,}71828-1)\approx42{,}957$. Valeur moyenne :
$\tfrac15\times42{,}957\approx\mathbf{8{,}59}$ milliers de têtes.

\noindent\textbf{Exercice 4.}
1) $z_C-z_A=2+4i$, $z_B-z_A=4$, donc $\dfrac{z_C-z_A}{z_B-z_A}=\dfrac{2+4i}{4}=\dfrac12+i$.
Ce quotient n'est ni réel ni imaginaire pur et a pour module
$\big|\tfrac12+i\big|=\sqrt{\tfrac14+1}=\tfrac{\sqrt5}{2}\neq1$ : le triangle $ABC$ est
quelconque (ni isocèle en $A$, ni rectangle en $A$).
2) $S:\ z'=az+b$ avec $S(A)=B$ et $S(B)=C$. On a
$a=\dfrac{z_C-z_B}{z_B-z_A}=\dfrac{(4+6i)-(6+2i)}{4}=\dfrac{-2+4i}{4}=-\tfrac12+i$.
Puis $b=z_B-a\,z_A=(6+2i)-(-\tfrac12+i)(2+2i)$. Or $(-\tfrac12+i)(2+2i)=-1-i+2i+2i^2=-1+i-2=-3+i$,
d'où $b=(6+2i)-(-3+i)=9+i$. Donc $z'=(-\tfrac12+i)z+9+i$.
3) Rapport $=|a|=\big|-\tfrac12+i\big|=\sqrt{\tfrac14+1}=\dfrac{\sqrt5}{2}$ ;
angle $=\arg(a)=\arg(-\tfrac12+i)=\pi-\arctan(2)\approx2{,}034$ rad (soit $\approx117{,}5°$).
4) Centre : $\omega=a\omega+b \iff \omega(1-a)=b \iff \omega=\dfrac{b}{1-a}$.
$1-a=1-(-\tfrac12+i)=\tfrac32-i$. Donc
$\omega=\dfrac{9+i}{\tfrac32-i}=\dfrac{(9+i)(\tfrac32+i)}{(\tfrac32)^2+1}
=\dfrac{\tfrac{27}{2}+9i+\tfrac32 i+i^2}{\tfrac94+1}
=\dfrac{\tfrac{27}{2}-1+\tfrac{21}{2}i}{\tfrac{13}{4}}
=\dfrac{\tfrac{25}{2}+\tfrac{21}{2}i}{\tfrac{13}{4}}=\dfrac{2(25+21i)}{13}=\dfrac{50+42i}{13}$.
Ainsi $\omega=\dfrac{50}{13}+\dfrac{42}{13}i\approx3{,}85+3{,}23\,i$.

\smallskip
\noindent\textbf{Partie B — Situation.}
\emph{Tâche 1.} Soient $x$ la largeur (deux côtés perpendiculaires au hangar) et $y$ la
longueur (côté parallèle au hangar, non doublé). Le grillage donne $2x+y=120$, soit
$y=120-2x$. Aire $A(x)=x(120-2x)=120x-2x^2$ ; $A'(x)=120-4x=0$ en $x=30$. Maximum en
$x=\SI{30}{\metre}$, $y=\SI{60}{\metre}$ ; aire maximale $A(30)=30\times60=\SI{1800}{\metre\squared}$.

\emph{Tâche 2.} Le cheptel double quand $N(t)=2\times400=800$, soit
$\mathrm{e}^{0{,}12t}=2 \iff t=\dfrac{\ln 2}{0{,}12}\approx5{,}78$. Le cheptel aura doublé
au bout de la \textbf{6\textsuperscript{e} année}.

\emph{Tâche 3.} Le nombre de bêtes réagissant mal suit une loi binomiale de paramètres
$n=5$ et $p=0{,}06$. $P(\text{au moins une})=1-(0{,}94)^5=1-0{,}7339\approx\mathbf{0{,}266}$,
soit environ \SI{26,6}{\percent}.
